\newpage
\subsection{\SurroundName}

Two clusterings are examined for \SurroundName, obtained at layer 8b with $K=4$ and at layer 7b with $K=6$.
At the deeper layer, besides the start of the episode, the clustering separates configurations by the region of the level occupied by the agent's trail, from a small trail confined to the bottom-left, to one running along the top after arriving from the left, to a dense trail that has grown from the top-left toward the bottom and is closing in on the final enclosure between the two players.
At the shallower layer, as for \PongName at $K=6$, this is an instance of over-segmentation: only four distinct trail configurations are visible among the six clusters, two of which are each duplicated.
The duplicated clusters leave \gls{cu}, \gls{as}, and \gls{acu} essentially unchanged from the $K=4$ clustering at layer 8b.

\begin{table}[H]
    \centering
    \begin{tblr}{
        colspec={|c|c|c|c|c|},
        hlines,
        vlines,
        row{1} = {font=\bfseries, halign=c, valign=m}
    }
        Layer & $K$ & \gls{cu}       & \gls{as}       & \gls{acu}      \\
        8b    & 4   & 0.67 $\pm$0.02 & 0.71 $\pm$0.02 & 0.38 $\pm$0.02 \\
        7b    & 6   & 0.67 $\pm$0.04 & 0.71 $\pm$0.04 & 0.38 $\pm$0.04 \\
    \end{tblr}
    \caption{\gls{cu}, \gls{as}, and \gls{acu} for the two clusterings analyzed in the \SurroundName environment.}
    \label{tab:qualitative_surround_metrics}
\end{table}

\begin{figure}[H]
    \qualitativeClusterFigureSetup
    \qualitativeClusterTwo
        {figures/clustering_agent_representation/pictures/surround/qualitative/cluster_1/image_1.png}
        {figures/clustering_agent_representation/pictures/surround/qualitative/cluster_1/image_2.png}
        {Start of the episode.}
    \qualitativeClusterTwo
        {figures/clustering_agent_representation/pictures/surround/qualitative/cluster_0/image_1.png}
        {figures/clustering_agent_representation/pictures/surround/qualitative/cluster_0/image_2.png}
        {The agent's trail remains confined to the bottom-left area.}
    \qualitativeClusterTwo
        {figures/clustering_agent_representation/pictures/surround/qualitative/cluster_2/image_1.png}
        {figures/clustering_agent_representation/pictures/surround/qualitative/cluster_2/image_2.png}
        {The agent moves along the top of the level, arriving from the left.}
    \qualitativeClusterTwo
        {figures/clustering_agent_representation/pictures/surround/qualitative/cluster_3/image_1.png}
        {figures/clustering_agent_representation/pictures/surround/qualitative/cluster_3/image_2.png}
        {The trail grows from the top-left toward the bottom, closing in on the final enclosure.}
    \caption{Clusters identified by \gls{car} in the \SurroundName environment (layer 8b, $K=4$).}
    \label{fig:qualitative_surround}
\end{figure}


\begin{figure}[H]
    \qualitativeClusterFigureSetup
    \qualitativeClusterTwo
        {figures/clustering_agent_representation/pictures/surround/qualitative_alt/cluster_0/image_1.png}
        {figures/clustering_agent_representation/pictures/surround/qualitative_alt/cluster_0/image_2.png}
        {E}
    \qualitativeClusterTwo
        {figures/clustering_agent_representation/pictures/surround/qualitative_alt/cluster_1/image_1.png}
        {figures/clustering_agent_representation/pictures/surround/qualitative_alt/cluster_1/image_2.png}
        {F}
    \qualitativeClusterTwo
        {figures/clustering_agent_representation/pictures/surround/qualitative_alt/cluster_2/image_1.png}
        {figures/clustering_agent_representation/pictures/surround/qualitative_alt/cluster_2/image_2.png}
        {G}
    \qualitativeClusterTwo
        {figures/clustering_agent_representation/pictures/surround/qualitative_alt/cluster_3/image_1.png}
        {figures/clustering_agent_representation/pictures/surround/qualitative_alt/cluster_3/image_2.png}
        {H}
    \qualitativeClusterTwo
        {figures/clustering_agent_representation/pictures/surround/qualitative_alt/cluster_4/image_1.png}
        {figures/clustering_agent_representation/pictures/surround/qualitative_alt/cluster_4/image_2.png}
        {I}
    \qualitativeClusterTwo
        {figures/clustering_agent_representation/pictures/surround/qualitative_alt/cluster_5/image_1.png}
        {figures/clustering_agent_representation/pictures/surround/qualitative_alt/cluster_5/image_2.png}
        {J}
    \caption{Clusters E--J identified by \gls{car} in the \SurroundName environment (layer 7b, $K=6$).}
    \label{fig:qualitative_surround_alt}
\end{figure}

